%% capitoli/cap7_de_iure_condendo.tex
%% In testa al file, prima di scrivere, incollare l'estratto di
%% output/F4-indice.md relativo a questo capitolo (tesi sostenuta,
%% fonti, lunghezza stimata): serve a controllare gli scostamenti.
%% Convenzione delle etichette: sez:cap7:<nome>, tab:cap7:<nome>.

% Ogni criticita' va classificata in una delle tre categorie: la proposta
% che ignora la classificazione non e' utilizzabile.
% Dimensionare il capitolo come un capitolo ordinario: e' cio' che
% distingue una tesi di specializzazione da una rassegna.

\chapter{Prospettive de iure condendo}
\label{cap:cap7}

\section{Le criticità risolvibili in sede di decreto legislativo entro i principi di delega}
\label{sez:cap7:delega}


\section{Le criticità che richiedono un intervento legislativo ulteriore}
\label{sez:cap7:nuovalegge}


\section{Quanto resta alla prassi della Corte e all'interpretazione}
\label{sez:cap7:prassi}


\section{Proposte, in ordine di priorità}
\label{sez:cap7:proposte}


